\documentclass[12pt,a4paper]{article}
\usepackage[margin=1in]{geometry}
\usepackage{amsmath,amssymb}
\usepackage{xcolor}
\usepackage{enumitem}
\usepackage{hyperref}
\usepackage{array}
\usepackage{booktabs}

\hypersetup{colorlinks=true,linkcolor=blue,urlcolor=blue}

\title{\textbf{ERD -- Entity Relationship Diagram}\\
\large Course Tracker \& Exam Notes}
\author{}
\date{}

\begin{document}
\maketitle
\tableofcontents
\newpage

\section{ERD Overview}

\subsection{What is an ERD?}
An Entity Relationship Diagram (ERD) is a visual representation of the data in a system. It shows:
\begin{itemize}
\item Entities
\item Attributes
\item Relationships between entities
\item Cardinality and participation
\item Keys such as primary keys
\end{itemize}

\textbf{Main purpose:} ERD helps design the database structure before implementing tables.

\subsection{Why use an ERD?}
\begin{itemize}
\item Understand system requirements.
\item Identify what data needs to be stored.
\item Identify relationships between different types of data.
\item Reduce redundancy and poor database design.
\item Provide a blueprint for converting requirements into database tables.
\end{itemize}

\subsection{Exam Definition}
\textbf{ERD:} A diagram that represents entities, their attributes, and the relationships between entities in a database system.

\section{Entity}

\subsection{What is an Entity?}
An entity is a real-world object or concept about which the system stores information.

Examples:
\begin{itemize}
\item Student
\item Employee
\item Department
\item Course
\item Customer
\item Product
\item Order
\end{itemize}

In a traditional ER diagram, an entity is represented using a \textbf{rectangle}.

\subsection{Entity Set}
An entity set is a collection of similar entities.

Example:
\[
Student(Student\_ID, Name, Age, Email)
\]

Each individual student is an entity, while the collection of students forms the entity set.

\section{Attributes}

\subsection{What is an Attribute?}
An attribute describes a property or characteristic of an entity.

For example:
\[
Student(Student\_ID, Name, Age, Email)
\]

Here, Student\_ID, Name, Age, and Email are attributes.

In a traditional ER diagram, attributes are represented using \textbf{ovals}.

\section{Types of Attributes}

\subsection{Simple Attribute}
An attribute that cannot be meaningfully divided into smaller attributes.

Example:
\[
Age
\]

\subsection{Composite Attribute}
An attribute that can be divided into smaller components.

Example:
\[
Name = FirstName + LastName
\]

Another example:
\[
Address = Street + City + State + PIN
\]

\subsection{Single-Valued Attribute}
An attribute that has one value for each entity.

Example: Student\_ID.

\subsection{Multi-Valued Attribute}
An attribute that may contain multiple values for one entity.

Example: PhoneNumber, because a student may have multiple phone numbers.

In a traditional ER diagram, a multi-valued attribute is represented using a \textbf{double oval}.

\subsection{Derived Attribute}
An attribute whose value can be calculated from another attribute.

Example:
\[
Age = CurrentYear - BirthYear
\]

In a traditional ER diagram, a derived attribute is represented using a \textbf{dashed oval}.

\section{Keys}

\subsection{Primary Key}
A primary key uniquely identifies each entity.

Example:
\[
Student(Student\_ID, Name, Age)
\]

Student\_ID can be the primary key.

Important properties:
\begin{itemize}
\item Unique
\item Cannot be NULL
\item Identifies an entity uniquely
\end{itemize}

In a traditional ER diagram, the key attribute is usually \textbf{underlined}.

\[
\underline{Student\_ID}
\]

\subsection{Candidate Key}
A candidate key is an attribute or set of attributes that can uniquely identify an entity. One candidate key is selected as the primary key.

\subsection{Composite Key}
A composite key consists of more than one attribute.

Example:
\[
Enrollment(Student\_ID, Course\_ID)
\]

The combination $(Student\_ID, Course\_ID)$ can uniquely identify an enrollment.

\section{Relationships}

\subsection{What is a Relationship?}
A relationship describes an association between entities.

Example:
\[
Student \longrightarrow Enrolls \longrightarrow Course
\]

In a traditional ER diagram, a relationship is represented using a \textbf{diamond}.

\subsection{Relationship Degree}

\subsubsection{Unary Relationship}
A relationship involving one entity type.

Example:
\[
Employee \longrightarrow Manages \longrightarrow Employee
\]

An employee manages another employee.

\subsubsection{Binary Relationship}
A relationship involving two entity types.

Example:
\[
Student \longrightarrow Enrolls \longrightarrow Course
\]

This is the most common type.

\subsubsection{Ternary Relationship}
A relationship involving three entity types.

Example: a Supplier supplies a Product to a Project, where Supplier, Product, and Project participate in one relationship.

\section{Cardinality}

Cardinality describes how many instances of one entity can be associated with instances of another entity.

\subsection{One-to-One (1:1)}
One entity is related to at most one entity on the other side.

Example:
\[
Person \longleftrightarrow Passport
\]

\subsection{One-to-Many (1:N)}
One entity can be related to many entities.

Example:
\[
Department \longrightarrow Employee
\]

One department can have many employees.

\subsection{Many-to-One (N:1)}
Many entities are related to one entity.

Example:
\[
Employee \longrightarrow Department
\]

Many employees can belong to one department.

\subsection{Many-to-Many (M:N)}
Many entities on one side can be related to many entities on the other side.

Example:
\[
Student \longleftrightarrow Course
\]

A student can enroll in many courses, and a course can have many students.

\section{Participation Constraints}

\subsection{Total Participation}
Every entity in the entity set must participate in the relationship.

Example: every employee must belong to a department.

\subsection{Partial Participation}
Some entities may participate in the relationship, but participation is not mandatory.

Example: not every employee necessarily manages another employee.

\section{Weak Entity}

A weak entity is an entity that cannot be uniquely identified using its own attributes alone.

It depends on another entity, called the \textbf{owner} or \textbf{strong entity}.

Example:
\[
Employee \longrightarrow Dependent
\]

A dependent may be identified using:
\[
Employee\_ID + DependentName
\]

\subsection{Strong vs Weak Entity}

\begin{center}
\begin{tabular}{p{0.42\textwidth} p{0.42\textwidth}}
\toprule
\textbf{Strong Entity} & \textbf{Weak Entity}\\
\midrule
Has its own identifying key & Depends on another entity for identification\\
Can exist independently & Existence depends on an owner entity\\
Example: Employee & Example: Dependent\\
\bottomrule
\end{tabular}
\end{center}

\section{ERD Symbols -- Quick Revision}

\begin{center}
\begin{tabular}{p{0.32\textwidth} p{0.48\textwidth}}
\toprule
\textbf{Component} & \textbf{Traditional Symbol}\\
\midrule
Entity & Rectangle\\
Relationship & Diamond\\
Attribute & Oval\\
Key Attribute & Underlined attribute\\
Multi-valued Attribute & Double oval\\
Derived Attribute & Dashed oval\\
Weak Entity & Double rectangle\\
Identifying Relationship & Double diamond\\
\bottomrule
\end{tabular}
\end{center}

\section{Example: Student-Course System}

Suppose:
\begin{itemize}
\item A student has a Student ID, name, and email.
\item A course has a Course ID and course name.
\item A student can enroll in multiple courses.
\item A course can contain multiple students.
\item The enrollment has an enrollment date.
\end{itemize}

Entities:
\[
Student(Student\_ID, Name, Email)
\]
\[
Course(Course\_ID, CourseName)
\]

Relationship:
\[
Student \longrightarrow Enrolls \longleftarrow Course
\]

Cardinality:
\[
Student \quad M:N \quad Course
\]

The relationship can have the attribute:
\[
EnrollmentDate
\]

This attribute belongs naturally to the enrollment relationship.

\section{Converting ERD to Relational Tables}

\subsection{Entity becomes a Table}
\[
Student(Student\_ID, Name, Email)
\]

becomes:
\begin{verbatim}
STUDENT
--------------------------
Student_ID   PRIMARY KEY
Name
Email
\end{verbatim}

\subsection{1:N Relationship}
Suppose:
\[
Department \quad 1:N \quad Employee
\]

The primary key of the 1-side is generally placed as a foreign key in the N-side.

\begin{verbatim}
DEPARTMENT
----------------
Department_ID  PK
DepartmentName

EMPLOYEE
----------------
Employee_ID    PK
EmployeeName
Department_ID  FK
\end{verbatim}

\subsection{M:N Relationship}
An M:N relationship is normally converted into a separate table.

\begin{verbatim}
STUDENT
----------------
Student_ID  PK
Name

COURSE
----------------
Course_ID   PK
CourseName

ENROLLMENT
----------------
Student_ID  FK
Course_ID   FK
EnrollmentDate
\end{verbatim}

The combination $(Student\_ID, Course\_ID)$ can be used as a composite primary key.

\section{ERD vs UML Class Diagram}

\begin{center}
\begin{tabular}{p{0.42\textwidth} p{0.42\textwidth}}
\toprule
\textbf{ERD} & \textbf{UML Class Diagram}\\
\midrule
Database/data modeling & Software/object modeling\\
Focuses on data & Focuses on objects/classes\\
Entities and relationships & Classes, methods, relationships\\
Attributes & Attributes and operations\\
Used for database design & Used for software design\\
\bottomrule
\end{tabular}
\end{center}

\[
ERD \rightarrow Database
\]
\[
UML \rightarrow Software/System Design
\]

\section{How to Draw an ERD from Requirements}

\begin{enumerate}
\item Read the requirements carefully.
\item Identify nouns -- they are potential entities.
\item Identify properties -- these become attributes.
\item Identify unique identifiers -- these become keys.
\item Identify verbs -- they often represent relationships.
\item Determine relationship cardinality.
\item Determine participation constraints if required.
\item Identify weak entities and identifying relationships.
\item Identify relationship attributes.
\item Review the complete ERD.
\end{enumerate}

\section{Common Exam Questions}

\subsection*{What is an ERD?}
A graphical representation of entities, their attributes, and relationships in a database system.

\subsection*{What is an entity?}
A real-world object or concept about which data is stored.

\subsection*{What is an attribute?}
A property that describes an entity.

\subsection*{What is cardinality?}
It specifies how many instances of one entity can be associated with instances of another entity.

\subsection*{What is a weak entity?}
An entity that cannot be uniquely identified by its own attributes alone and depends on an owner entity.

\subsection*{What is a composite attribute?}
An attribute that can be divided into smaller meaningful attributes.

\subsection*{What is a derived attribute?}
An attribute whose value is calculated from other stored data.

\section{Exam Checklist}

\begin{itemize}[label=$\square$]
\item ERD definition and purpose
\item Entity and entity set
\item Attributes
\item Simple, composite, single-valued, multi-valued, and derived attributes
\item Primary, candidate, and composite keys
\item Relationships
\item Unary, binary, and ternary relationships
\item 1:1, 1:N, N:1, and M:N cardinality
\item Total and partial participation
\item Strong and weak entities
\item ERD symbols
\item Converting ERD to relational tables
\item ERD vs UML class diagram
\item Drawing an ERD from requirements
\end{itemize}

\section{One-Page Memory Sheet}

\begin{center}
\textbf{\Large ERD = Entities + Attributes + Relationships}
\end{center}

\begin{verbatim}
ENTITY
  |
  +-- ATTRIBUTES
  |
  +-- KEY

ENTITY ---- RELATIONSHIP ---- ENTITY
              |
           CARDINALITY
          /    |     \
        1:1   1:N    M:N
\end{verbatim}

\textbf{Remember:}

\begin{verbatim}
Rectangle        -> Entity
Oval             -> Attribute
Diamond          -> Relationship
Underlined       -> Key attribute
Double Oval      -> Multi-valued attribute
Dashed Oval      -> Derived attribute
Double Rectangle -> Weak entity

1:1 -> One to One
1:N -> One to Many
M:N -> Many to Many

ERD -> Database design
UML -> Software design
\end{verbatim}

\end{document}
